\documentclass[]{article}
\usepackage{lmodern}
\usepackage{amssymb,amsmath}
\usepackage{ifxetex,ifluatex}
\usepackage{fixltx2e} % provides \textsubscript
\ifnum 0\ifxetex 1\fi\ifluatex 1\fi=0 % if pdftex
  \usepackage[T1]{fontenc}
  \usepackage[utf8]{inputenc}
\else % if luatex or xelatex
  \ifxetex
    \usepackage{mathspec}
  \else
    \usepackage{fontspec}
  \fi
  \defaultfontfeatures{Ligatures=TeX,Scale=MatchLowercase}
\fi
% use upquote if available, for straight quotes in verbatim environments
\IfFileExists{upquote.sty}{\usepackage{upquote}}{}
% use microtype if available
\IfFileExists{microtype.sty}{%
\usepackage{microtype}
\UseMicrotypeSet[protrusion]{basicmath} % disable protrusion for tt fonts
}{}
\usepackage[left=1.1in, right=1.1in, top=1.1in, bottom=1.1in]{geometry}
\usepackage{hyperref}
\hypersetup{unicode=true,
            pdftitle={Stat 435 Lecture Notes 4b},
            pdfauthor={Xiongzhi Chen; Washington State University},
            pdfborder={0 0 0},
            breaklinks=true}
\urlstyle{same}  % don't use monospace font for urls
\usepackage{color}
\usepackage{fancyvrb}
\newcommand{\VerbBar}{|}
\newcommand{\VERB}{\Verb[commandchars=\\\{\}]}
\DefineVerbatimEnvironment{Highlighting}{Verbatim}{commandchars=\\\{\}}
% Add ',fontsize=\small' for more characters per line
\usepackage{framed}
\definecolor{shadecolor}{RGB}{248,248,248}
\newenvironment{Shaded}{\begin{snugshade}}{\end{snugshade}}
\newcommand{\KeywordTok}[1]{\textcolor[rgb]{0.13,0.29,0.53}{\textbf{{#1}}}}
\newcommand{\DataTypeTok}[1]{\textcolor[rgb]{0.13,0.29,0.53}{{#1}}}
\newcommand{\DecValTok}[1]{\textcolor[rgb]{0.00,0.00,0.81}{{#1}}}
\newcommand{\BaseNTok}[1]{\textcolor[rgb]{0.00,0.00,0.81}{{#1}}}
\newcommand{\FloatTok}[1]{\textcolor[rgb]{0.00,0.00,0.81}{{#1}}}
\newcommand{\ConstantTok}[1]{\textcolor[rgb]{0.00,0.00,0.00}{{#1}}}
\newcommand{\CharTok}[1]{\textcolor[rgb]{0.31,0.60,0.02}{{#1}}}
\newcommand{\SpecialCharTok}[1]{\textcolor[rgb]{0.00,0.00,0.00}{{#1}}}
\newcommand{\StringTok}[1]{\textcolor[rgb]{0.31,0.60,0.02}{{#1}}}
\newcommand{\VerbatimStringTok}[1]{\textcolor[rgb]{0.31,0.60,0.02}{{#1}}}
\newcommand{\SpecialStringTok}[1]{\textcolor[rgb]{0.31,0.60,0.02}{{#1}}}
\newcommand{\ImportTok}[1]{{#1}}
\newcommand{\CommentTok}[1]{\textcolor[rgb]{0.56,0.35,0.01}{\textit{{#1}}}}
\newcommand{\DocumentationTok}[1]{\textcolor[rgb]{0.56,0.35,0.01}{\textbf{\textit{{#1}}}}}
\newcommand{\AnnotationTok}[1]{\textcolor[rgb]{0.56,0.35,0.01}{\textbf{\textit{{#1}}}}}
\newcommand{\CommentVarTok}[1]{\textcolor[rgb]{0.56,0.35,0.01}{\textbf{\textit{{#1}}}}}
\newcommand{\OtherTok}[1]{\textcolor[rgb]{0.56,0.35,0.01}{{#1}}}
\newcommand{\FunctionTok}[1]{\textcolor[rgb]{0.00,0.00,0.00}{{#1}}}
\newcommand{\VariableTok}[1]{\textcolor[rgb]{0.00,0.00,0.00}{{#1}}}
\newcommand{\ControlFlowTok}[1]{\textcolor[rgb]{0.13,0.29,0.53}{\textbf{{#1}}}}
\newcommand{\OperatorTok}[1]{\textcolor[rgb]{0.81,0.36,0.00}{\textbf{{#1}}}}
\newcommand{\BuiltInTok}[1]{{#1}}
\newcommand{\ExtensionTok}[1]{{#1}}
\newcommand{\PreprocessorTok}[1]{\textcolor[rgb]{0.56,0.35,0.01}{\textit{{#1}}}}
\newcommand{\AttributeTok}[1]{\textcolor[rgb]{0.77,0.63,0.00}{{#1}}}
\newcommand{\RegionMarkerTok}[1]{{#1}}
\newcommand{\InformationTok}[1]{\textcolor[rgb]{0.56,0.35,0.01}{\textbf{\textit{{#1}}}}}
\newcommand{\WarningTok}[1]{\textcolor[rgb]{0.56,0.35,0.01}{\textbf{\textit{{#1}}}}}
\newcommand{\AlertTok}[1]{\textcolor[rgb]{0.94,0.16,0.16}{{#1}}}
\newcommand{\ErrorTok}[1]{\textcolor[rgb]{0.64,0.00,0.00}{\textbf{{#1}}}}
\newcommand{\NormalTok}[1]{{#1}}
\usepackage{graphicx,grffile}
\makeatletter
\def\maxwidth{\ifdim\Gin@nat@width>\linewidth\linewidth\else\Gin@nat@width\fi}
\def\maxheight{\ifdim\Gin@nat@height>\textheight\textheight\else\Gin@nat@height\fi}
\makeatother
% Scale images if necessary, so that they will not overflow the page
% margins by default, and it is still possible to overwrite the defaults
% using explicit options in \includegraphics[width, height, ...]{}
\setkeys{Gin}{width=\maxwidth,height=\maxheight,keepaspectratio}
\IfFileExists{parskip.sty}{%
\usepackage{parskip}
}{% else
\setlength{\parindent}{0pt}
\setlength{\parskip}{6pt plus 2pt minus 1pt}
}
\setlength{\emergencystretch}{3em}  % prevent overfull lines
\providecommand{\tightlist}{%
  \setlength{\itemsep}{0pt}\setlength{\parskip}{0pt}}
\setcounter{secnumdepth}{0}
% Redefines (sub)paragraphs to behave more like sections
\ifx\paragraph\undefined\else
\let\oldparagraph\paragraph
\renewcommand{\paragraph}[1]{\oldparagraph{#1}\mbox{}}
\fi
\ifx\subparagraph\undefined\else
\let\oldsubparagraph\subparagraph
\renewcommand{\subparagraph}[1]{\oldsubparagraph{#1}\mbox{}}
\fi

%%% Use protect on footnotes to avoid problems with footnotes in titles
\let\rmarkdownfootnote\footnote%
\def\footnote{\protect\rmarkdownfootnote}

%%% Change title format to be more compact
\usepackage{titling}

% Create subtitle command for use in maketitle
\newcommand{\subtitle}[1]{
  \posttitle{
    \begin{center}\large#1\end{center}
    }
}

\setlength{\droptitle}{-2em}

  \title{Stat 435 Lecture Notes 4b}
    \pretitle{\vspace{\droptitle}\centering\huge}
  \posttitle{\par}
    \author{Xiongzhi Chen \\ Washington State University}
    \preauthor{\centering\large\emph}
  \postauthor{\par}
    \date{}
    \predate{}\postdate{}
  
\usepackage{bbm}
\usepackage{amssymb}
\usepackage{amsmath}
\usepackage{graphicx,float}

\begin{document}
\maketitle

{
\setcounter{tocdepth}{2}
\tableofcontents
}
\section{}\label{section}

\section{Ridge regression: basics}\label{ridge-regression-basics}

\subsection{Overview}\label{overview}

Ridge regression

\begin{itemize}
\tightlist
\item
  applies regardless of magnitues of sample size \(n\) and number of
  predictors \(p\)
\item
  shrinks estimated coefficients compared to \emph{least squares
  estimate (LSE)}
\item
  is able to produce LSE
\item
  presents a way to estimate coefficients when \(n < p+1\)
\end{itemize}

\emph{Ridge regression often produces a biased estimate with smaller
variance than LSE, and it is a bias-variance trade-off technique.}

\subsection{Settings}\label{settings}

\begin{itemize}
\item
  Model:
  \(Y=\beta_0+\beta_1 X_1 + \beta_2 X_2 + \ldots + \beta_p X_p + \varepsilon\)
\item
  Observations: \((y_i,x_{i1},x_{i2},\ldots,x_{ip}),i=1,\ldots,n\),
  where \(y_i\) is the \(i\)th observation for \(Y\) and \(x_{ij}\) that
  for \(X_j\)
\item
  Estimate
  \(\hat{\boldsymbol{\beta}}=(\hat{\beta}_1,\ldots,\hat{\beta}_p)\) of
  \({\boldsymbol{\beta}}=({\beta}_1,\ldots,{\beta}_p)\), and
  \(\hat{\beta}_0\) of \({\beta}_0\)
\item
  Fitted model:
  \(\hat{y}_i=\hat{\beta}_0+\hat{\beta}_1 x_{i1} + \hat{\beta}_2 x_{i2} + \ldots + \hat{\beta}_p x_{ip}\)
\item
  Residuals: \(e_i=y_i - \hat{y}_i\)
\end{itemize}

\subsection{Optimization}\label{optimization}

\begin{itemize}
\tightlist
\item
  The ridge estimate
  \(\hat{\boldsymbol{\beta}}^R_{\lambda}=(\hat{\beta}_1,\ldots,\hat{\beta}_p)\)
  is the \({\boldsymbol{\beta}}=({\beta}_1,\ldots,{\beta}_p)\) that
  minimizes \[
  L_2(\beta_0,\boldsymbol{\beta},\lambda)= \sum_{i=1}^n (y_i - \hat{y}_i)^2 + \lambda \sum_{i=1}^p \beta_i^2
  \]
\item
  \(\sqrt{\sum_{i=1}^p \beta_i^2}\) is written as
  \(\Vert\boldsymbol{\beta}\Vert_2\), i.e., the \(l_2\)-norm of
  \(\boldsymbol{\beta}\)\\
\item
  \(\sum_{i=1}^n (y_i - \hat{y}_i)^2\) is just the RSS
\item
  \(\lambda \sum_{i=1}^p \beta_i^2\) is called a \emph{regularizer} or
  \emph{penalty}
\item
  \(L_2(\beta_0,\boldsymbol{\beta},\lambda)\) is referred to as an
  \emph{\(l_2\)-regularized RSS}
\item
  no penalty on \(\beta_0\)
\end{itemize}

\subsection{Intercept term}\label{intercept-term}

\begin{itemize}
\tightlist
\item
  Recall \(\mathbf{x}_i=(x_{i1},x_{i2},\ldots,x_{ip}),i=1,\ldots,n\),
  where \(x_{ij}\) is the \(i\)th observation for \(X_j\)
\item
  Let \(\mathbf{X}\) be the \(n \times p\) matrix whose \(i\)th row is
  \(\mathbf{x}_i\)
\item
  \(\mathbf{X}\) is called a \emph{design matrix} (when factors as
  predictors are properly coded)
\end{itemize}

If the variables \(X_i,i=1,\ldots,p\) are centered, i.e., the columns of
\(\mathbf{X}\), to have mean zero before ridge regression is performed,
then the estimated intercept will be
\[\hat{\beta}_0=\bar{y}=\sum_{i=1}^n y_i/n\]

\subsection{Solution path}\label{solution-path}

Recall the objective function \[
L_2(\beta_0,\boldsymbol{\beta},\lambda)= \sum_{i=1}^n (y_i - \hat{y}_i)^2 + \lambda \sum_{i=1}^p \beta_i^2
\] and its solution
\(\hat{\boldsymbol{\beta}}^R_{\lambda}=(\hat{\beta}_1,\ldots,\hat{\beta}_p)\)

\begin{itemize}
\tightlist
\item
  no penalty on estimated \(\beta_0\)
\item
  \(\lambda=0\) gives the LSE, and \(\lambda=\infty\) forces
  \(\hat{\boldsymbol{\beta}}^R_{\lambda}=0\)
\item
  some \(0 < \lambda < \infty\) strikes a balance between LSE and
  \(\hat{\boldsymbol{\beta}}^R_{\lambda}=0\)
\end{itemize}

\emph{Note:} The ridge solution has explicit reprentation.

\subsection{Bias-variance trade-off}\label{bias-variance-trade-off}

Ridge regression works best in situations where the least squares
estimates have high variance:

\begin{itemize}
\tightlist
\item
  \(\lambda=0\): ridge estimate is the LSE and is unbiased
\item
  \(\lambda >0\): bias increases and variance decreases usually
\item
  \(\lambda=\infty\): estimated coefficients are all zero
\end{itemize}

\subsection{Bias-variance trade-off}\label{bias-variance-trade-off-1}

\(p=45\) predictors and \(n=50\) observations; all \(\beta_j \ne 0\)

\begin{center}\includegraphics[width=0.85\linewidth]{fig6_5} \end{center}

\subsection{Select tuning parameter}\label{select-tuning-parameter}

The optimal value \(\lambda^{\ast}\) of the tuning parameter \(\lambda\)
is ofen determined by \(k\)-fold cross-validation:

\begin{itemize}
\tightlist
\item
  Pick a sequence of \(s\) values for \(\lambda\) as
  \(\lambda_1, \lambda_2, \ldots, \lambda_s\)
\item
  For each \(\lambda_l\), apply \(k\)-fold cross-validation to estimate
  the test error of the corresponding model
\item
  Set \(\lambda^{\ast}\) as the value among the \(s\) values of
  \(\lambda\) for which the estimated test error is the smallest
\end{itemize}

\section{Ridge regression: special
case}\label{ridge-regression-special-case}

\subsection{Special model}\label{special-model}

\begin{itemize}
\tightlist
\item
  Assume the design matrix \(\mathbf{X}=\mathbf{I}_p\), i.e., the
  identity matrix
\item
  Consider a linear regression model without an intercept, i.e., forcing
  \(\beta_0=0\)
\end{itemize}

If \(n=p\) and \(\mathbf{X}=\mathbf{I}_p\), then we have a very special
model: \[y_j = \beta_j+\varepsilon_j,j=1,\ldots,p\]

\emph{Note:} \(\mathbf{X}=\mathbf{I}_p\) is referred to as an
\emph{orthogonal design}

\subsection{LSE and ridge estimate}\label{lse-and-ridge-estimate}

For the special model \[y_j = \beta_j+\varepsilon_j,j=1,\ldots,p,\]

\begin{itemize}
\tightlist
\item
  the LSE (i.e., least squares estimate) is \[\hat{\beta}_j=y_j\]
\item
  the \(l_2\)-regularized RSS becomes \[
  L_2(\beta_{0},\boldsymbol{\beta},\lambda)= \sum_{j=1}^p (y_i - \beta_j)^2 + \lambda \sum_{i=1}^p \beta_i^2
  \] and the ridge estimate is
  \[\hat{\beta}_{j,\lambda}^R=y_j/(1+\lambda)\]
\end{itemize}

\subsection{Three estimators}\label{three-estimators}

\begin{center}\includegraphics[width=0.85\linewidth]{fig6_10} \end{center}

\section{Ridge regression:
application}\label{ridge-regression-application}

\subsection{Scale equivariance}\label{scale-equivariance}

\begin{itemize}
\tightlist
\item
  Recall the design matrix \(\mathbf{X}\), whose \((i,j)\)th entry
  \(x_{ij}\) is the \(i\)th observation of predictor \(X_j\)
\item
  Ridge regression estimates depend on the scale of predictors
\item
  Before applying ridge regression, it is best to standardize predictors
  as follows \[
  \tilde{x}_{ij}= \frac{x_{ij}}{\sqrt{\frac{1}{n}\sum_{i=1}^n (x_{ij}-\bar{x}_j)^2}}
  \] where \(\bar{x}_j=n^{-1}\sum_{i=1}^n x_{ij}\)
\item
  The resulting ridge gression estimates are called \emph{standardized
  estimates}
\end{itemize}

\subsection{Illustration}\label{illustration}

Modelling the \texttt{Credit} data set

\begin{itemize}
\tightlist
\item
  Response \texttt{Balance}
\item
  Predictors \texttt{Income}, \texttt{Limit}, \texttt{Rating} and
  \texttt{Student}
\end{itemize}

\subsection{Illustration}\label{illustration-1}

\begin{center}\includegraphics[width=0.85\linewidth]{fig6_4} \end{center}

\subsection{Inference}\label{inference}

\begin{itemize}
\tightlist
\item
  The method of \emph{bias correction} can be used to conduct inference
  on ridge estimates
\item
  Canonical assumptions involve Gaussian random errors
\item
  Asymptotic theory on testing coefficients is based on Gaussian
  limiting distributions
\item
  P-values from testing can be obtained
\end{itemize}

\emph{Note:} ``High-Dimensional Inference: Confidence Intervals,
p-Values and R-Software hdi'' by Ruben Dezeure, Peter Buhlmann, Lukas
Meier and Nicolai Meinshausen

\subsection{Inference}\label{inference-1}

\begin{itemize}
\tightlist
\item
  Tesing \(H_{j0}:\beta_j =0\) versus \(H_{j1}:\beta_j \ne 0\)
\item
  P-values for testing if the coefficient of each of \texttt{Income},
  \texttt{Limit}, \texttt{Rating} and \texttt{Student}:
\end{itemize}

\begin{verbatim}
       Income         Limit        Rating    StudentYes 
1.183510e-236  1.655880e-15  6.850389e-19 1.808087e-132 
\end{verbatim}

\section{LASSO: basics}\label{lasso-basics}

\subsection{Settings}\label{settings-1}

\begin{itemize}
\item
  Model:
  \(Y=\beta_0+\beta_1 X_1 + \beta_2 X_2 + \ldots + \beta_p X_p + \varepsilon\)
\item
  Observations: \((y_i,x_{i1},x_{i2},\ldots,x_{ip}),i=1,\ldots,n\),
  where \(y_i\) is the \(i\)th observation for \(Y\) and \(x_{ij}\) that
  for \(X_j\)
\item
  Estimate
  \(\hat{\boldsymbol{\beta}}=(\hat{\beta}_1,\ldots,\hat{\beta}_p)\) of
  \({\boldsymbol{\beta}}=({\beta}_1,\ldots,{\beta}_p)\), and
  \(\hat{\beta}_0\) of \({\beta}_0\)
\item
  Fitted model:
  \(\hat{y}_i=\hat{\beta}_0+\hat{\beta}_1 x_{i1} + \hat{\beta}_2 x_{i2} + \ldots + \hat{\beta}_p x_{ip}\)
\item
  Residuals: \(e_i=y_i - \hat{y}_i\)
\end{itemize}

\subsection{Optimization}\label{optimization-1}

\begin{itemize}
\tightlist
\item
  The LASSO estimate
  \(\hat{\boldsymbol{\beta}}^L_{\lambda}=(\hat{\beta}_1,\ldots,\hat{\beta}_p)\)
  is the \({\boldsymbol{\beta}}=({\beta}_1,\ldots,{\beta}_p)\) that
  minimizes \[
  L_1(\beta_0,\boldsymbol{\beta},\lambda)= \frac{1}{2}\sum_{i=1}^n (y_i - \hat{y}_i)^2 + \lambda \sum_{i=1}^p\vert \beta_i \vert
  \]
\item
  \(\sum_{i=1}^p \vert\beta_i\vert\) is written as
  \(\Vert\boldsymbol{\beta}\Vert_1\), i.e., the \(l_1\)-norm of
  \(\boldsymbol{\beta}\)\\
\item
  \(\sum_{i=1}^n (y_i - \hat{y}_i)^2\) is just the RSS
\item
  \(\lambda \sum_{i=1}^p \vert\beta_i\vert\) is called a
  \emph{regularizer} or \emph{penalty}
\item
  \(L_1(\beta_0,\boldsymbol{\beta},\lambda)\) is referred to as an
  \emph{\(l_1\)-regularized RSS}
\item
  no penalty on \(\beta_0\)
\end{itemize}

\subsection{Solution path}\label{solution-path-1}

Recall the objective function \[
L_1(\beta_0,\boldsymbol{\beta},\lambda)= \frac{1}{2}\sum_{i=1}^n (y_i - \hat{y}_i)^2 + \lambda \sum_{i=1}^p \vert \beta_i \vert
\] and its solution
\(\hat{\boldsymbol{\beta}}^L_{\lambda}=(\hat{\beta}_1,\ldots,\hat{\beta}_p)\)

\begin{itemize}
\tightlist
\item
  no penalty on estimated \(\beta_0\)
\item
  \(\lambda=0\) gives the LSE, and \(\lambda=\infty\) forces
  \(\hat{\boldsymbol{\beta}}^L_{\lambda}=0\)
\item
  some \(0 < \lambda < \infty\) strikes a balance between LSE and
  \(\hat{\boldsymbol{\beta}}^L_{\lambda}=0\)
\item
  \emph{some \(\hat{\beta}_j\)'s,\(j=1,\ldots,p,\) can be exactly zero}
\end{itemize}

\subsection{Bias-variance trade-off}\label{bias-variance-trade-off-2}

LASSO works best in situations where some coefficeints are extactly
zero:

\begin{itemize}
\tightlist
\item
  \(\lambda=0\): LASSO estimate is the LSE and is unbiased
\item
  \(\lambda >0\): bias increases and variance decreases usually, and
  some estimated coefficients are zero
\item
  \(\lambda=\infty\): estimated coefficients are all zero
\end{itemize}

\subsection{Select tuning parameter}\label{select-tuning-parameter-1}

The optimal value \(\lambda^{\ast}\) of the tuning parameter \(\lambda\)
is ofen determined by \(k\)-fold cross-validation:

\begin{itemize}
\tightlist
\item
  Pick a sequence of \(s\) values for \(\lambda\) as
  \(\lambda_1, \lambda_2, \ldots, \lambda_s\)
\item
  For each \(\lambda_l\), apply \(k\)-fold cross-validation to estimate
  the test error of the corresponding model
\item
  Set \(\lambda^{\ast}\) as the value among \(\lambda\) as
  \(\lambda_1, \lambda_2, \ldots, \lambda_s\) for which the estimated
  test error is the smallest
\end{itemize}

\section{LASSO estimate: special
case}\label{lasso-estimate-special-case}

\subsection{Special model}\label{special-model-1}

\begin{itemize}
\tightlist
\item
  Recall \(\mathbf{x}_i=(x_{i1},x_{i2},\ldots,x_{ip}),i=1,\ldots,n\),
  where \(x_{ij}\) is the \(i\)th observation for \(X_j\)
\item
  Let \(\mathbf{X}\) be the \(n \times p\) matrix whose \(i\)th row is
  \(\mathbf{x}_i\)
\item
  Consider a linear regression model without an intercept, i.e., forcing
  \(\beta_0=0\)
\end{itemize}

If \(n=p\) and \(\mathbf{X}=\mathbf{I}_p\), then we have a very special
model: \[y_j = \beta_j+\varepsilon_j,j=1,\ldots,p\]

\subsection{LSE and LASSO estimate}\label{lse-and-lasso-estimate}

For the special model \[y_j = \beta_j+\varepsilon_j,j=1,\ldots,p,\]

\begin{itemize}
\tightlist
\item
  the LSE (i.e., least squares estimate) is \[\hat{\beta}_j=y_j\]
\item
  the \(l_1\)-regularized RSS becomes \[
  L_1(\beta_{0},\boldsymbol{\beta},\lambda)= \sum_{j=1}^p (y_i - \beta_j)^2 + \lambda \sum_{i=1}^p \vert\beta_i\vert
  \]
\end{itemize}

\subsection{LSE and LASSO estimate}\label{lse-and-lasso-estimate-1}

The LASSO estimate, more complicated than LSE and ridge estimate, is:

\[\hat{\beta}_{j,\lambda}^L = \left\{
\begin{array}
{lll}
0 & \text{if} & \vert y_j \vert \le \lambda \\
y_j -\lambda & \text{if} & y_j > \lambda \\
y_j +\lambda & \text{if} & y_j < -\lambda
\end{array}\right.
\] \emph{Note:} compare the above with LSE \(\hat{\beta}_j=y_j\) and
ridge estimate \[\hat{\beta}_{j,\lambda}^R=y_j/(1+\lambda)\]

\subsection{Three estimators}\label{three-estimators-1}

\begin{center}\includegraphics[width=0.85\linewidth]{fig6_10} \end{center}

\section{LASSO: application}\label{lasso-application}

\subsection{Illustration}\label{illustration-2}

Modelling the \texttt{Credit} data set

\begin{itemize}
\tightlist
\item
  Response \texttt{Balance}
\item
  Predictors \texttt{Income}, \texttt{Limit}, \texttt{Rating} and
  \texttt{Student}
\end{itemize}

\subsection{Illustration}\label{illustration-3}

\begin{center}\includegraphics[width=0.85\linewidth]{fig6_6} \end{center}

\subsection{Illustration}\label{illustration-4}

\(p=45\) predictors and \(n=50\) observations; 43 \(\beta_j\)'s are 0

\begin{center}\includegraphics[width=0.85\linewidth]{fig6_13} \end{center}

\subsection{Inference}\label{inference-2}

\begin{itemize}
\tightlist
\item
  The method of \emph{bias correction} can be used to conduct inference
  on LASSO estimates
\item
  Canonical assumptions involve Gaussian random errors
\item
  Asymptotic theory on testing coefficients is based on Gaussian
  limiting distributions
\item
  P-values from testing can be obtained
\end{itemize}

\emph{Note:} ``High-Dimensional Inference: Confidence Intervals,
p-Values and R-Software hdi'' by Ruben Dezeure, Peter Buhlmann, Lukas
Meier and Nicolai Meinshausen

\subsection{Inference}\label{inference-3}

\begin{itemize}
\tightlist
\item
  Tesing \(H_{j0}:\beta_j =0\) versus \(H_{j1}:\beta_j \ne 0\)
\item
  P-values for testing if the coefficient of each of \texttt{Income},
  \texttt{Limit}, \texttt{Rating} and \texttt{Student}:
\end{itemize}

\begin{verbatim}
       Income         Limit        Rating    StudentYes 
2.160475e-255 1.039480e-101 1.923041e-133 1.896622e-128 
\end{verbatim}

\section{Ridge and LASSO estimates}\label{ridge-and-lasso-estimates}

\subsection{Three optimizations}\label{three-optimizations}

\begin{center}\includegraphics[width=0.85\linewidth]{fig6_7} \end{center}

\subsection{Three estimators}\label{three-estimators-2}

\begin{center}\includegraphics[width=0.85\linewidth]{fig6_10} \end{center}

\subsection{Ridge and LASSO
estimates}\label{ridge-and-lasso-estimates-1}

\(p=45\) predictors and \(n=50\) observations; all \(\beta_j \ne 0\)

\begin{center}\includegraphics[width=0.85\linewidth]{fig6_8} \end{center}

\subsection{Ridge and LASSO
estimates}\label{ridge-and-lasso-estimates-2}

\(p=45\) predictors and \(n=50\) observations; 43 \(\beta_j\)'s are 0

\begin{center}\includegraphics[width=0.85\linewidth]{fig6_9} \end{center}

\subsection{Ridge and LASSO
regressions}\label{ridge-and-lasso-regressions}

\begin{itemize}
\tightlist
\item
  Ridge regression performs better when the response is a function of
  many predictors, all with coefficients of roughly equal size
\item
  LASSO performs better when a relatively small number of predictors
  have substantial coefficients and the remaining predictors have very
  small or zero coefficients
\item
  Both ridge and LASSO regression yield a reduction in variance at the
  expense of a small increase in biases, when the LSE have excessively
  high variance
\item
  LASSO performs variable selection while ridge regression does not
\end{itemize}

\subsection{License and session
Information}\label{license-and-session-information}

\href{http://math.wsu.edu/faculty/xchen/stat412/LICENSE.html}{License}

\begin{Shaded}
\begin{Highlighting}[]
\NormalTok{>}\StringTok{ }\KeywordTok{sessionInfo}\NormalTok{()}
\NormalTok{R version }\FloatTok{3.5.0} \NormalTok{(}\DecValTok{2018-04-23}\NormalTok{)}
\NormalTok{Platform:}\StringTok{ }\NormalTok{x86_64-w64-mingw32/}\KeywordTok{x64} \NormalTok{(}\DecValTok{64}\NormalTok{-bit)}
\NormalTok{Running under:}\StringTok{ }\NormalTok{Windows }\DecValTok{10} \KeywordTok{x64} \NormalTok{(build }\DecValTok{19045}\NormalTok{)}

\NormalTok{Matrix products:}\StringTok{ }\NormalTok{default}

\NormalTok{locale:}
\NormalTok{[}\DecValTok{1}\NormalTok{] LC_COLLATE=English_United States}\FloatTok{.1252} 
\NormalTok{[}\DecValTok{2}\NormalTok{] LC_CTYPE=English_United States}\FloatTok{.1252}   
\NormalTok{[}\DecValTok{3}\NormalTok{] LC_MONETARY=English_United States}\FloatTok{.1252}
\NormalTok{[}\DecValTok{4}\NormalTok{] LC_NUMERIC=C                          }
\NormalTok{[}\DecValTok{5}\NormalTok{] LC_TIME=English_United States}\FloatTok{.1252}    

\NormalTok{attached base packages:}
\NormalTok{[}\DecValTok{1}\NormalTok{] stats     graphics  grDevices utils     datasets  methods  }
\NormalTok{[}\DecValTok{7}\NormalTok{] base     }

\NormalTok{other attached packages:}
\NormalTok{[}\DecValTok{1}\NormalTok{] knitr_1}\FloatTok{.21}

\NormalTok{loaded via a }\KeywordTok{namespace} \NormalTok{(and not attached):}
\StringTok{ }\NormalTok{[}\DecValTok{1}\NormalTok{] compiler_3}\FloatTok{.5.0}  \NormalTok{magrittr_1}\FloatTok{.5}    \NormalTok{tools_3}\FloatTok{.5.0}    
 \NormalTok{[}\DecValTok{4}\NormalTok{] htmltools_0}\FloatTok{.3.6} \NormalTok{yaml_2}\FloatTok{.2.0}      \NormalTok{Rcpp_1}\FloatTok{.0.3}     
 \NormalTok{[}\DecValTok{7}\NormalTok{] stringi_1}\FloatTok{.2.4}   \NormalTok{rmarkdown_1}\FloatTok{.11}  \NormalTok{stringr_1}\FloatTok{.3.1}  
\NormalTok{[}\DecValTok{10}\NormalTok{] xfun_0}\FloatTok{.4}        \NormalTok{digest_0}\FloatTok{.6.18}   \NormalTok{evaluate_0}\FloatTok{.12}  
\end{Highlighting}
\end{Shaded}


\end{document}
